\chapter{Ultracold Gases and quantum mixtures}
\begin{flushright}{\slshape{~}
	%It is only in very recent times that it has begun to dawn upon\\ the mind of man that there is another world only one degree\\ less complex and wonderful than the stellar universe \\
	%-- the world of molecules and atoms \\ \medskip
	%\textnormal{--- Albert A. Michelson \cite{Michelson1896ar}} %\defcitealias{knuth:1974}{Donald E. Knuth}\citetalias{knuth:1974} 

}
\end{flushright}

\section{Why ultracold atoms?}

The year $1925$ was a particularly important year for the field of quantum physics. Through the contributions of Heisenberg, Born, Schr\"odinger and many others, decisive steps forward were made towards a deeper theoretical footing of quantum physics. In particular, Heisenberg's ``Umdeutung'' paper \cite{Heisenberg1925uqu} drew a sharp distinction from the deterministic approach of classical mechanics. 

In the same year, Wolfgang Pauli attempted to provide a theoretical basis for   Bohr's empirical closed-shell atomic model, and at the same time to understand the Zeeman effect. He postulated the \graffito{A more formal expression of the exclusion principle concerns the anti-symmetry properties of many-body wave function}\textit{Pauli exclusion principle} \cite{Pauli1925udz}, which states that only a single electron in an atomic orbital can occupy a specific state, that is, no two electrons can share the same set of quantum numbers.

A year after Pauli, exactly a hundred years ago, Fermi \cite{Fermi1926zqd}  and Dirac \cite{Dirac1926ott} independently extended Pauli's results to the case of a non-interacting gas consisting of indistinguishable particles, obeying the exclusion principle. These results laid the groundwork for what is now known as Fermi-Dirac statistics, which describes the thermodynamic behavior, at the quantum level, of particles with half-integer spin, known as fermions. The Pauli principle shapes most of the world around us, as fermions constitute the fundamental building blocks of matter and the exclusion principle dictates how subnuclear, nuclear, and atomic particles organize themselves.

The Fermi-Dirac framework is complementary to the one proposed, around the same time, first by Bose \cite{Bose1924pgu} and then by Einstein \cite{Einstein1924qde}. They considered the situation in which non-interacting particles are fully indistinguishable and are allowed to be all in the same quantum state. In this case particles with integer spin, called bosons, follow the Bose-Einstein statistics. A consequence of this statistic is that, in the limit of zero temperature, a phase transition occurs and bosonic particles macroscopically occupy the lowest available energy state, forming a \ac{BEC}, a macroscopic quantum state described by a single wave function.

The deep connection between the spin of a particle and the statistics it obeys was formulated by Pauli in the spin-statistic theorem \cite{Pauli1940tcb}. This theorem implies that a quantum theory is not only probabilistic, but also has an inherently statistical mechanical character. While these results are valid for all physical objects for which a quantum description is required, for a long time the interplay between statistics and many-body quantum effects remained difficult to observe directly.

The first hint of Bose-Einstein statistics in nature came from the observation of \textit{superfluidity} in liquid helium. Below the critical temperature of $T_c=\SI{2.17}{K}$, the flow of liquid helium-$4$, a bosonic isotope of helium, was observed to become frictionless \cite{Kapitza1938vol, Allen1938fpi} and this was interpreted by Fritz London in terms of the formation of a \ac{BEC}, which allowed the particles to flow without viscosity \cite{London1938tlp}. However, this intuition was only qualitative, Einstein's theory for condensation was developed for a non-interacting gas and liquid helium was only observed to become superfluid in an interacting system. For a long time, further developments were halted by the lack of an experimental platform to investigate this phenomenon.

A similar statistical mechanical reasoning was later applied to \textit{superconductivity}. Superconductivity is the phenomenon whereby current flows through a conductor without electrical resistance, under certain conditions, and it was first observed in $1911$ by Onnes \cite{Onnes1912few}, who was investigating the conductivity of mercury cooled to temperatures of a few kelvins using liquid helium. A microscopic theoretical explanation remained elusive until the development of the \ac{BCS} theory \cite{Bardeen1957tos}. In the proposed theory, electrons with opposite spin, despite being fermions, would form a composite boson known as a \textit{Cooper pair} and condense into the lowest energy state, forming a charged superfluid in which electron pairs flow without resistance. While this theory successfully explained experimental observation of superconductivity at low temperatures, as in the case of superfluidity, direct experimental observation of the underlying microscopic mechanisms remained out of reach, and highly controllable systems with tunable interactions were needed in order to study the microscopic behavior of these phenomena. 

Atomic gases emerged as such a platform, and considerable progress towards a controlled experimental study of quantum statistical phenomena was possible thanks to the development of laser cooling techniques for atomic gases in the $1980$s and early $1990$s. The availability of cooling techniques was of crucial importance, as in order to  experimentally observe quantum effects in an atomic cloud, ultra-low temperatures, of the order of hundreds of nanokelvin or below, are needed. At these temperatures quantum mechanical phenomena become predominant and manifest at the macroscopic scale. Owing to the wave-particle duality of matter, each atom can be assigned a thermal de Broglie wavelength, expressed as
\begin{equation}
	\lambda_{\textrm th}=\frac{h}{\sqrt{2\pi mk_BT}}\,,
\end{equation}
with $h$ the Planck constant, $k_B$ the Boltzmann constant, $m$ the mass of atoms and $T$ the temperature. When $\lambda_{\textrm th}$ is of the order of the interparticle distance, i.e.~$\lambda_{\textrm th}\sim n^{-1/3}$, where 
$n$ is the number density of the gas, the quantum statistical nature of the atoms plays a fundamental role.


However, the temperatures and densities achieved with laser cooling were not enough to reach quantum degeneracy\footnote{It took almost $20$ years after the first realization for a \ac{BEC} to be achieved through direct laser cooling \cite{Stellmer2013lct}.}, and the final step was made possible by evaporative cooling of the trapped gas \cite{Hess1986eco, Ketterle1996eco}.
The first milestone was reached in $1995$, with the observation of \ac{BEC} in a dilute ultracold quantum gas of rubidium, at JILA \cite{Anderson1995oob}. Soon after, \ac{BEC} was also achieved in a gas of sodium, at MIT \cite{Davis1995bec}. A few years later, in $1999$, mirroring the historical development of the respective statistics, the first \ac{DFG} was realized with ultracold potassium \cite{DeMarco1999oof}, also at JILA, and shortly after with lithium, in many other laboratories around the world  \cite{Truscott2001oof, Schreck2001qbe, Hadzibabic2002tsm, Jochim2004PhD}.

With both quantum statistics experimentally realized with ultracold atoms, many new exciting opportunities arose to study quantum systems in a controlled and reproducible manner. Thanks to their great tunability, ultracold atoms provide a highly controllable platform for testing theoretical predictions. Among the numerous possibilities, in this thesis we will focus on the behavior of systems composed of more than one species of quantum gas, known as a quantum mixtures.




\section{Quantum mixtures}%two is better than one

The first realizations of mixtures of two degenerate quantum gases were achieved already shortly after the realizations of the first \ac{BEC}s. These first few attempts focused on mixtures of two different spin states of the same atomic species \cite{Myatt1997pot, Stenger1998sdi} as they were the most straight forward and accessible combination to realize. 

While the access to mixture in bosonic system is interesting in order to increase the complexity of the system under study and to allow the exploration of more complex and interesting phases, in fermionic gases operating with a mixture is (almost) mandatory. Generally, when a fermionic gas composed of a single component is cooled to low temperatures, the interaction freezes out, as at low energies only short-range scattering is relevant. Owing to the Pauli exclusion principle, the wave function overlap between identical fermions has to be zero and thus \textit{s}-wave interactions are forbidden. This is a twofold problem, not only does the interesting physics require interactions, but the final cooling step to reach degeneracy, evaporative cooling, relies on elastic collisions in the atomic gas to rethermalize the cloud. 

The first degenerate Fermi gases were indeed realized in mixtures, either with a mixture of two spin states of the alkali isotope of interest, as in the case of $^{40}$K \cite{DeMarco1999oof} and many of the $^6$Li realizations, see e.g.\,Ref.~\cite{Granade2002aop}, or by \textit{sympathetic cooling} of fermionic atoms with another better-behaved species, for which evaporative cooling is available. In this case, the atomic species of choice can be of two types. On one hand, it is possible to use a different isotope of the atoms under investigation to allow interactions. This is the situation, e.g., of the first realizations of degenerate $^6$Li, that was achieved by mixing it with a cloud of its bosonic brother $^7$Li \cite{Truscott2001oof, Schreck2001qbe}. On the other hand, the alternative is to use a different atomic species with better know scattering properties. This technique was not only employed to reach quantum degeneracy with fermionic gases \cite{Hadzibabic2002tsm}, but also with bosonic species which proved to be more challenging to cool than others. This is the case, for example, of bosonic potassium, which was condensed for the first time by sympathetic cooling with rubidium \cite{Modugno2001bec}.

These three approaches to cooling serves also to exemplify the three kinds of mixtures realizable with quantum gases. The experimentally easiest and most natural approach is to mix two (or potentially more) spin states of the same atomic species in a homonuclear mixture. This configuration requires minimal differences in the experimental setup, compared to a single-species experiment, as in most cases it does not require laser cooling of additional components. For bosonic mixtures, the spin mixture can be realized with the gas already in the ultracold regime, while in most of the fermionic cases, a spin mixture is required to perform evaporative cooling. Somewhat more demanding experimentally are isotope mixtures. In this case some parts of the experimental setup are shared between the species involved, e.g.~the atomic source, and optical components can be shared between the isotopes. The most experimentally demanding configuration for quantum mixtures is reached with a mixture of two different atomic species. In this case both atomic species require separate sources and the optical transitions involved can have considerably different properties, requiring dedicated lasers for each species. Each additional step in complexity, however, has the possibility of enabling the access of different ranges of physical phenomena, see Refs.~\cite{Baroni2024qmo,Varenna2022book} for some extensive but not exhaustive examples, and thus the choice of species and platform is in itself one of the most critical experimental decision.

Spin mixtures are bound to work with Bose-Bose or Fermi-Fermi combinations since they originate from a single atomic species. \graffito{In a broad sense heteronuclear mixture can consist of different isotopes or different atomic species}\textit{Heteronuclear} mixtures instead do not suffer this kind of limitation and thus they can also feature Bose-Fermi mixtures. As we will see in the next section however, heteronuclear systems are limited by the quantum statistics of the atomic species and isotopes available in nature and by the feasibility to laser cool them. If suitable isotopes of the chosen atomic species are available, in heteronuclear mixtures experiments it is possible to change quantum statistics of the components if needed, with limited experimental downsides.



%add some sentences on spin mixture

%reread cos review, nice to present some example of the broad range of mixtures before moving to mass imbalanced and fermions

%Bose-Bose mixtures, whether spin-based or heteronuclear, have proven to be an extremely successful platform to study numerous phenomena. From spinor \ac{BEC} to false vacuum decay to creation of ground state molecules and other very nice examples...  see e.g.~Ref.~ for a (non-exhaustive) compendium of the many possibility of physical phenomena available with quantum mixtures.


% Summarizing properly the results obtained in the broad field of quantum mixtures is no easy task. Many books and review have been written on the topic and to pick some precise example among all the different mixtures and results would be in part to pick favourite

Our focus in this thesis is dedicated towards the characterization of the interaction properties of an unusual heteronuclear quantum mixture, mixing two fermionic isotopes. In the following section, we provide an overview of the mass-imbalance heteronuclear mixture and present the motivations for the choice of atomic species employed in our experiment. 

% 


\section{Why D\lowercase{y}-K?} 

Heteronuclear mixtures find application across a wide range of contexts, from serving as a starting point for the assembly of ultracold ground-state molecules \cite{Bigagli2024oob}, to providing a platform for the study of complex interactions and exotic many-body phases.\albi{	add some more ref}

The objective of our experiment is to study how the presence of mass imbalance could affect the pairing mechanism in fermionic superfluids. While the \ac{BCS} theory successfully explains the microscopic behavior of superconductivity in so-called \graffito{Superconductors are also categorized according to the 
}\textit{conventional} systems, superconductors with critical temperatures of several tens of kelvin or higher are not explained within the same framework. These \textit{unconventional} superconductors, observed for the first time in $1986$ \cite{Bednorz1986pht}, are of particular interest for numerous potential industrial applications and the open theoretical questions they pose.

Mass-balanced, homonuclear spin mixtures of ultracold Fermi gases proved to be a formidable platform to better understand conventional superconductors, providing a highly controllable experimental platform on which to test theoretical predictions. In particular, the precise interaction tuning, enabled by access to magnetic Feshbach resonances \cite{Chin2010fri}, allowed for a detailed study of the pairing mechanism across the \ac{BEC}-\ac{BCS} crossover \cite{Parish2014tbb}, that is, the smooth transition from tightly bound molecular pairs to loosely correlated Cooper pairs as the interaction strength and sign is varied.

Similarly, mass-imbalanced Fermi mixtures hold the potential to deepen our understanding of unconventional superconductors. Exotic pairing mechanisms, different from the standard Cooper pair formation, are considered to be possible explanations for these \graffito{Unconventional superconductivity often occurs at higher critical temperatures than for conventional systems}\textit{high-temperature} superconductors \albi{cite}. Theoretical calculations predict that the difference in mass between the two components of a pair could allow ultracold gases to access exotic phases of matter analogous to those present in unconventional superconductors. These phases are also predicted for spin mixtures; however, the temperature conditions required for their realization are not experimentally attainable in such systems, while for heteronuclear mixtures they are predicted to occur under more realistic conditions. \albi{cite}


Moreover, similar unconventional pairing mechanisms are also considered to be present in other strongly interacting fermionic systems, such as neutron stars and the primordial quark-gluon plasma present in the early universe \albi{cite}. Regardless of mass imbalance, strongly interacting fermionic systems for which the characteristic interaction length scale diverges are said to be in the \textit{unitary regime} and the physics described in this framework has a universal character. This universal behavior observed in ultracold atoms allows to deepen our understanding of vastly different systems that are in the same regime, like it is predicted to occur in neutron stars and in the early universe. \albi{cite} %some stuff of Martin for example}


%when scattering length diverges, the physics in fermionic systems becomes  universal, as the scattering length becomes the only relevant length (and energy) scale


%, but experimental realization of these pairing mechanisms in experiments with a high degree of control is still lacking. 

%few possible candidates to explore, and even less explored


The first natural candidate to study a mass-imbalanced fermionic system is a mixture of lithium and potassium, as the only two alkali species with stable\,\footnote{$^{40}$K is technically radioactive, with a half-life of \SI{1.248e9}{y} \cite{Kondev2021tne}, only one order of magnitude shorter than the age of the universe \cite{Planck2020p2r}, which is plenty of time for a cold-atoms experiment.} fermionic isotopes. Both species individually proved themselves to be excellent platforms to study strongly interacting fermionic mixtures, and pairing and superfluidity across the \ac{BEC}-\ac{BCS} crossover, see e.g.~Refs.~\cite{Jochim2003bec, Zwierlein2003oob,Zwierlein2004cop,Zwierlein2005vas,Zwierlein2006fsw, Greiner2003eoa, Regal2004oor} for some of the first results of the community on the topic, and the proceedings of the $2006$ E.~Fermi school in Varenna \cite{Varenna2006book} devoted to the topic. 


The initial experimental realizations of the Li-K Fermi-Fermi mixture appeared in $2008$ \cite{Wille2008eau,Tiecke2009fri,Taglieber2008qdt}. Over the following two decades, this mixture proved to be an extremely fruitful platform, with results ranging from the study of hydrodynamic behavior \cite{Trenkwalder2011heo}, to the properties of Feshbach \cite{Voigt2009uhf,Jag2016lof} and ground-state molecules \cite{He2024eco} and, in particular in the Innsbruck experiment, to the study of impurity physics \cite{Kohstall2012mac, Cetina2016umb, Baroni2024mib}. Considerable theoretical work was also dedicated to the role of the Li-K mass ratio in the \ac{BEC}-\ac{BCS} crossover \cite{Yi2006, Mackie2005, Gandolfi2024bcs}.

The Feshbach resonances available in the Li-K mixture however proved to be excessively narrow and prone to two- and three-body losses \cite{Naik2011fri,Jag2016lof}, which hindered the study of pair superfluidity. This limitation, combined with the increasing number of new atomic species with fermionic isotopes brought to quantum degeneracy over the last decades, such as  ytterbium \cite{Fukuhara2007dfg}, strontium \cite{DeSalvo2010dfg}, dysprosium \cite{Lu2012qdd}, erbium \cite{Aikawa2014rfd}, and chromium \cite{Naylor2015cdf}, motivated the exploration of alternative combinations.

In order to choose a suitable mixture for our goals, the mass ratio between the two components has to be carefully chosen. If the ratio is too small, the effect of the mass imbalance would be negligible, while for mass ratios $M/m>12.33$, where $M$ ($m$) is the mass of the heavier (lighter) specie, a mixture is expected to suffer from large three-body recombination losses \cite{Petrov2005dmi}. A further constraint regards the the electronic structure of the atomic species involved. Alkaline-earth-like atoms present a closed shell in their valence energy levels. Atoms with a closed shell, like strontium and ytterbium, have been observed not to support broad Feshbach resonances \cite{Green2020fri, Barbe2018oof,Mukherjee2022fra}, and thus are not suitable for the goal of studying a strongly interacting system, further narrowing the range of possible combinations. 
  
As an additional constraint, the choice was made to keep one of the two species as an alkali, i.e.~to have either lithium or potassium, to avoid excessively increasing the experimental complexity. With this limitation in mind the choice was narrowed down to the combinations of Li-Cr, Dy-K, and Er-K. At the time when our experiment was started (in $2015$), fermionic chromium was not yet well characterized as it had only recently been brought to degeneracy, and lanthanide atoms with a large dipolar magnetic moment, like Er and Dy, were the subject of significant interest,  thanks to their long-range interactions \cite{Chomaz2022dpa}.

The extremely dense interaction spectrum of dysprosium and erbium, which has been observed to be chaotic \cite{Maier2015eoc, Maier2015buf}, makes it unfeasible to predict the possible interaction spectrum of a mixture involving one of the two species. Moreover, fermionic isotopes have a finite nuclear spin, which provides additional splitting to the electronic states compared to bosonic isotopes, making theoretical prediction of the interaction spectrum of the mixture too cumbersome. 

Therefore the final choice for our experiment was to use potassium-$40$ combined with dysprosium, which features two fermionic isotopes, $^{161}$Dy and $^{163}$Dy, as opposed to a single one of erbium, effectively doubling our opportunities of finding suitable interspecies Feshbach resonances. The history of the Dy-K experiment has thus been largely exploratory, directed towards the preparation for the first time of a doubly degenerate mixture composed of a lanthanide and an alkali, achieved in $2018$ and reported in Ref.~\cite{Ravensbergen2018poa}, identifying suitable scattering resonances and finally approaching the study of the \ac{BEC}-\ac{BCS} crossover with mass imbalance. 

In this thesis we focus on the last two objectives, building on the results already obtained by the previous member of the group. In the \hyperref[Chpt:Feshbach]{following chapter}, we give an overview of strongly interacting fermionic gases, relevant to both spin- and heteronuclear mixtures, useful to better understand the topics treated in this work. In Chapter~\ref{Chpt:Outline}, we outline the structure of the rest of this thesis,. by briefly summarizing each chapter.
 

% TODO, maybe relfect to the fact that other mixture have been realized since the incipit of this experiment, as it is now it seems a bit too much that this is the first Dy-K thesis, which is false.


%\footnote{This large dipole moment holds particular usefulness for an experiment working with fermionic isotopes as, opposite to what usually happens with non-magnetic atoms, intraspecies interaction are never completely frozen even at very low temperatures and quantum degeneracy can be reached in a single spin state, thanks to universal dipolar scattering \cite{Lu2012qdd,Aikawa2014rfd}.}





 
 
%In standard fermionic experiments, evaporative cooling is performed between two spin states of the same atomic species, in order to allow for the interaction necessary for rethermalization. 
